% 14-07(14쪽) — y=f(x) 의 그래프(x=-2 에서 끊기는 두 조각 · 왼쪽은 y=a-2 인 반직선, 오른쪽은 위로 볼록한 곡선).
% 스캔 화질이 낮아 TikZ 로 다시 그림. 원문의 빈 점(○)·채운 점(●)·점선과 눈금 a+2, a, a-2, -2 만 옮긴다.
\begin{tikzpicture}[x=0.55cm, y=0.55cm, line width=0.55pt]
  \draw[->, >=stealth, line width=0.7pt] (-3.15,0) -- (3.35,0) node[below=1pt] {$x$};
  \draw[->, >=stealth, line width=0.7pt] (0,-1.7) -- (0,3.5) node[left=1pt] {$y$};
  \draw[densely dashed, line width=0.4pt] (-2,1.95) -- (0,1.95);
  \draw[densely dashed, line width=0.4pt] (-2,1.15) -- (0,1.15);
  \draw[densely dashed, line width=0.4pt] (-2,0.35) -- (0,0.35);
  \draw[densely dashed, line width=0.4pt] (-2,1.95) -- (-2,0);
  \draw[line width=0.6pt] (-3.05,0.35) -- (-2,0.35);
  \draw[line width=0.6pt, smooth] plot coordinates {(-2,1.95) (-1.5,2.45) (-1,2.72) (-0.35,2.85) (0.3,2.78) (1.0,2.6) (1.6,2.1) (1.95,1.5) (2.2,0.8) (2.45,0) (2.7,-0.9) (2.85,-1.5)};
  \draw[fill=white] (-2,0.35) circle (2.2pt);
  \fill (-2,1.15) circle (2.2pt);
  \draw[fill=white] (-2,1.95) circle (2.2pt);
  \node[right=1pt, font=\small] at (0,1.95) {$a+2$};
  \node[right=1pt, font=\small] at (0,1.15) {$a$};
  \node[right=1pt, font=\small] at (0,0.35) {$a-2$};
  \node[below=1pt, font=\small] at (-2,0) {$-2$};
  \node[below left=-2pt and -4pt, font=\small] at (0,0) {$\pt{O}$};
  \node[anchor=west, font=\small] at (0.15,3.25) {$y=f(x)$};
\end{tikzpicture}
